The primary benchmark uses the named NumPy numerical-library operation or an integer recurrence converted to binary64 as its reference. In the supplementary elementary-function experiments, the MPFR library for multiple-precision floating-point computation supplies correctly rounded mathematical values \cite{mpfr}. \emph{Correct rounding} returns the representable number nearest the exact mathematical result, choosing the one with an even last significand bit in a tie. The \emph{significand} holds the significant binary digits and the exponent gives their scale. Normal binary64 numbers have 53 significant bits. A dyadic rational is an integer divided by a power of two: for example, $3/8$ is exactly $0.011$ in binary. Every finite binary64 number has this form.

A \emph{unit in the last place} (ULP) is one step in the last stored digit at a given scale \cite{goldberg1991}. We measure ULP distance by counting adjacent representable values, treating both signs of zero as equal. At $1$, for example, the next larger number is $1+2^{-52}$ and the next smaller number is $1-2^{-53}$. Both have distance one under this convention. The mean scaled absolute error averages each absolute prediction error divided by the larger of one and the absolute reference value. A parameter \emph{slot} stores one scalar even if it is zero. An \emph{edit} changes one stored scalar from its initialized value. An \emph{operation} is an executed arithmetic or activation step in the stated graph. A real-arithmetic proof assumes exact algebra; a binary64 proof specifies where rounding occurs.
